\section{Bakgrunn}
\label{sec:bakgrunn}

Helsetjenestens driftsorganisasjon for nødnett HF (HDO) leverer 
kommunikasjonsløsninger til norsk akuttmedisin og drifter 
videoløsninger for 109 akuttmedisinske kommunikasjonssentraler 
(AMK) og legevaktsentraler (LVS) over hele Norge 
\cite{hdo_kravspec}. Løsningene brukes daglig i akutte situasjoner 
der operatører sender en SMS-lenke til innringeren, som deretter 
kan sende videostrøm tilbake til sentralen. Lyd går via ordinær 
telefonsamtale; video brukes som beslutningsstøtte ved hendelser 
som hjertestans, trafikkulykker og fødsler.

HDO benytter to separate videomotorer: AntMedia Server og NHN 
Video (basert på Pexip). De to motorene har ulike API-er og 
integrasjonsmetoder, noe som medfører at ny funksjonalitet må 
implementeres to ganger, opplæring og vedlikehold skjer i to 
separate systemer, og terskelen for å bytte eller legge til 
videomotorer er høy. Siden systemet brukes i akuttmedisinsk 
sammenheng hele døgnet, stiller HDO strenge krav til stabilitet 
og tilgjengelighet.

For å løse dette ønsket HDO en mellomvare som skjuler forskjellene 
mellom videomotorene bak ett felles API. Dette prosjektet 
utviklet og leverte en slik mellomvare — VideoAPI — med integrasjon 
mot begge motorene.